This Design Document (DD) provides a detailed description of the technical and architectural design decisions for the Students\&Companies (S\&C) platform.
It builds upon the requirements outlined in the Requirements Analysis and Specification Document (RASD),
translating them into a comprehensive design framework for implementation, integration, and testing.
This document specifies the architecture, system components, and data flows that implement required functionalities while ensuring scalability, security, and compliance with non-functional requirements such as reliability and maintainability.

\section{Purpose}\label{sec:purpose}
The purpose of this document is to provide the foundational design and architecture required to implement the S\&C platform. It ensures that the platform:
\begin{itemize}
    \item Meets the functional and non-functional requirements outlined in the RASD.
    \item Supports secure and efficient interactions among students, companies, and universities.
    \item Integrates with external tools such as SMTP providers for notifications and statistical analysis tools for evaluating internship processes.
    \item Provides a robust foundation for testing, verification, and validation of the system.
\end{itemize}

\section{Scope}\label{sec:scope}
The Students\&Companies platform streamlines the internship process by connecting students with companies based on skills, experiences, and goals.
Its main features include:
\begin{itemize}
    \item Tools for students to search for internships and for companies to advertise opportunities.
    \item A recommendation engine to match students’ CVs with companies’ internship requirements.
    \item Management of interviews, questionnaires, and candidate evaluations.
    \item Feedback mechanisms for improving internship applications and postings.
    \item Monitoring tools for universities to track internship progress and handle complaints.
    \item Notifications for students and companies about relevant matches and updates.
\end{itemize}

The system incorporates a secure authentication mechanism for all user types (Company, Student, University)
and ensures the accuracy of user-provided data.
It employs a three-tier architecture to separate the user interface, application logic, and data storage,
allowing for scalability and maintainability.

\section{Glossary}\label{sec:glossary}

\subsection{Definitions}\label{subsec:definitions}
\begin{itemize}
    \item \textbf{Student:} A user* registered on the platform to search for internships and submit their CVs.
    \item \textbf{Company:} A user* registered on the platform to post internship opportunities and select candidates.
    \item \textbf{University:} A user* with the authority to monitor internships and handle complaints.
    \item \textbf{CV:} A document summarizing a student’s skills, experiences, and qualifications.
    \item \textbf{Internship:} A professional experience offered by a company to train students in specific domains.
    \item \textbf{Recommendation Engine:} A system that matches students and internships based on predefined criteria such as skills and experience.
    \item \textbf{Feedback:} Retrospective information provided by users* to evaluate internships and enhance the platform.
    \item \textbf{Complaint:} A report submitted by students, companies, or universities to address issues related to internships.
    \item \textbf{User:} Includes authenticated students and companies.
    \item \textbf{User*:} Includes authenticated students, companies and universities.
    \item \textbf{Guest:} Unauthenticated students, companies, or universities accessing limited features of the platform.
\end{itemize}

\subsection{Acronyms}\label{subsec:acronyms}
\begin{itemize}
    \item \textbf{S\&C}: Students\&Companies
    \item \textbf{RASD}: Requirements Analysis and Specification Document
    \item \textbf{DD}: Design Document
    \item \textbf{CV}: Curriculum Vitae
    \item \textbf{SMTP}: Simple Mail Transfer Protocol
    \item \textbf{GDPR}: General Data Protection Regulation
    \item \textbf{API}: Application Programming Interface
\end{itemize}

\subsection{Abbreviations}\label{subsec:abbreviations}
\begin{itemize}
    \item \textbf{e.g.}: Example (\textit{exempli gratia})
    \item \textbf{i.e.}: That is (\textit{id est})
    \item \textbf{etc.}: And so on (\textit{et cetera})
    \item \textbf{vs.}: Against (\textit{versus})
    \item \textbf{w.r.t.}: With respect to
\end{itemize}

\section{Revision History}\label{sec:revision-history}
\begin{table}[H]
    \centering
    \renewcommand{\arraystretch}{1.5}
    \begin{tabular}{|c|c|p{0.6\textwidth}|}
        \hline
        \rowcolor[HTML]{5C8A99}
        \textbf{Version} & \textbf{Date} & \textbf{Description} \\ \hline
        1.0              & 11/12/2024    & Initial version of the Design Document. \\ \hline
    \end{tabular}
    \caption{Revision History}
\end{table}

\section{Reference Documents}\label{sec:reference-documents}
\begin{itemize}
    \item Assignment RDD AY 2024-2025: Problem statement and project description.
    \item RASD: Requirements Analysis and Specification Document for the S\&C platform.
    \item IEEE Std 830-1998: IEEE Recommended Practice for Software Requirements Specifications.
    \item GDPR: General Data Protection Regulation for data privacy compliance.
    \item UML 2.0: Unified Modeling Language standards for diagrams and models.
    \item Software Engineering 2 Course Materials: Lectures and resources from the course.
\end{itemize}

\section{Document Structure}\label{sec:document-structure}
This document is structured as follows:
\begin{enumerate}
    \item \textbf{Introduction:} Overview of the document's purpose, scope, glossary, revision history, and references.
    \item \textbf{Architectural Design:} High-level system architecture, including component diagrams and deployment diagrams.
    \item \textbf{Detailed Design:} Comprehensive design specifications for each system component.
    \item \textbf{User Interface Design:} Description and wireframes for the platform’s user interface.
    \item \textbf{Requirements Traceability:} Mapping of requirements from the RASD to system components.
    \item \textbf{Implementation, Integration and Testing:} Description of the implementation approach and testing strategies.
    \item \textbf{Effort Spent:} Summary of contributions from each team member.
    \item \textbf{References:} List of all reference materials used during the project.
\end{enumerate}
